\section{Introduction} \label{sec:intro}

Autonomous systems, including unmanned aerial vehicles (UAVs), self-driving ground vehicles, and mobile robots, increasingly rely on networked communication protocols to coordinate missions, exchange telemetry, and receive operator commands.
The MAVLink protocol~\cite{mavlink2024}, widely adopted by platforms such as ArduPilot and PX4, transmits heartbeat signals, GPS coordinates, sensor readings, and command acknowledgements over UDP links that are often unencrypted and unauthenticated.
This architectural choice makes autonomous platforms inherently vulnerable to a range of cyber-physical attacks: message flooding, identity spoofing, replay injection, and command manipulation~\cite{mekdad2023uavsecurity, neshenko2019iot, cui2010embedded}.

Significant research effort has been devoted to individual attack vectors (e.g., GPS spoofing~\cite{shoukry2017safecontrol}, replay attacks~\cite{mo2009replay}, LiDAR spoofing~\cite{sun2020lidar}) and to specific defence mechanisms (e.g., federated IDS~\cite{nguyen2021fliot}, specification-based monitors~\cite{bartocci2018specmonitoring}, wireless IDS~\cite{mitchell2014wireless}).
However, these studies typically evaluate \emph{one} attack against \emph{one} defence in a bespoke experimental setup.
Consequently, results cannot be compared across papers: different simulators, traffic profiles, metric definitions, and mission scenarios make cross-study comparison unreliable~\cite{sommer2010closedworld}.
Most published experiments are also not reproducible because the underlying code, configurations, and datasets are rarely shared~\cite{collberg2016repeatability}.

\subsection{Research Gap}

We identify three fundamental gaps in the current state of the art:

\begin{enumerate}
    \item \textbf{No multi-attack, multi-defence orchestration.}
    Existing frameworks either inject faults without adversarial intent (DriveFI~\cite{jha2019drivefi}, software fault injection surveys~\cite{natella2016faultinjection}) or benchmark IDS classifiers on static datasets without live mission context (UAV-NIDD~\cite{hadi2025uavnidd}).
    No published tool runs multiple heterogeneous attacks against multiple defence strategies within the same live mission execution.

    \item \textbf{Lack of a centralised orchestrator.}
    Current evaluations require manual scripting to coordinate attack timing, defence activation, and metric collection.
    There is no reusable orchestration layer that schedules attacks (with deterministic or randomised timing), activates defences, synchronises logging, and computes standardised security metrics automatically.

    \item \textbf{Missing mission-level impact assessment.}
    Most IDS benchmarks report detection rate and false-positive rate but ignore the downstream effect on mission performance, such as trajectory deviation, completion delay, or system crash.
    For safety-critical autonomous platforms, understanding mission impact is as important as raw detection accuracy.
\end{enumerate}

\subsection{Research Goals and Contributions}

This paper presents the design, implementation, and evaluation of a \emph{Red-Team vs.\ Blue-Team Evaluation Framework} for autonomous platforms communicating over MAVLink.
The framework introduces a centralised orchestrator, a component absent from prior work, that automates the entire evaluation lifecycle: launching the simulated platform, scheduling parameterised attacks from a red-team catalogue, activating configurable blue-team defences, collecting synchronised logs, and computing standardised metrics.

Our contributions are:

\begin{enumerate}
    \item \textbf{Orchestrator-driven evaluation harness.}
    A reusable, open-source framework with abstract base classes for attacks (\texttt{BaseAttack}), defences (\texttt{BaseDefense}), detectors (\texttt{BaseDetector}), and platforms (\texttt{BasePlatform}).
    New modules can be added by implementing a single interface, without modifying the orchestrator.

    \item \textbf{Red-team attack catalogue.}
    Six MAVLink-specific attacks (heartbeat flood, ping flood, parameter-request flood, MITM identity spoof, replay-pattern injection, and command-injection burst) with configurable rate, duration, and randomised scheduling modes.

    \item \textbf{Formal problem definition and blue-team defence strategies.}
    A formal definition of the red-team vs.\ blue-team evaluation problem (platform, attacks, schedule, detector, defence, orchestrator, and metrics) together with three defence mechanisms: a Markov transition-probability filter, an Isolation Forest anomaly detector trained on 11 engineered features, and a rate-threshold baseline. All three are evaluated under identical conditions.

    \item \textbf{Systematic experimental evaluation with statistical analysis.}
    Four experiments comparing no-defence, static-attack with Markov defence, randomised-attack with Markov defence, and Isolation Forest adaptive defence, with multi-seed repetition (20 random seeds for randomised experiments), mean/standard-deviation reporting, statistical significance testing (Mann-Whitney U), and a baseline comparison against the rate-threshold detector.

    \item \textbf{Reproducibility.}
    All code, configurations, and generated datasets are publicly available.\footnote{\url{https://github.com/The-Christopher-Robin/autonomy-evaluation-harness}} Experiments are deterministically reproducible via command-line parameters and random seeds.
\end{enumerate}

The remainder of this paper is organised as follows.
Section~\ref{sec:related} surveys related work and positions our framework relative to existing tools.
Section~\ref{sec:framework} details the framework architecture, feature engineering, and defence algorithms.
Section~\ref{sec:experiments} presents experimental results.
Section~\ref{sec:discussion} discusses findings and limitations.
Section~\ref{sec:threats} addresses threats to validity.
Section~\ref{sec:conclusion} concludes with future directions.
